\chapter{Fabriquer une grille en JSON}

Ce chapitre est le plus court, et c'est significatif : \textbf{ajouter un plateau
ne demande pas une ligne de code.} Le chapitre suivant montre la voie inverse —
un plateau entièrement décrit en C++ — et explique quand la préférer.

\section{Pourquoi le plateau est une donnée}

\begin{nkcode}{REGLES\_COMPLETES\_v2.md:93-107}
Le plateau n'est pas une constante du code : c'est un descripteur serialisable,
editable a la souris dans le moteur de test.

BoardDesc {
  topology     : HEX_POINTY | HEX_FLAT | SQUARE_4 | SQUARE_8
  cells        : liste de coordonnees valides        // definit la forme reelle
  blocked      : liste de coordonnees bloquees
  starts       : [ { player_index, coord, level } ]  // positions initiales
  min_players  : entier
  max_players  : entier
}
\end{nkcode}

Les formes prévues pour Conqueror — rectangle, hexagone, plus, diamant — ne sont
pas dans le moteur : ce sont des fichiers. La raison est pratique. Le palier 0
doit répondre à « quelle taille de plateau ? quel voisinage ? quel avantage au
premier joueur ? » : trois questions qui demandent d'essayer vingt plateaux. Si
chaque essai coûte une recompilation, on en essaiera trois.

\section{Où déposer le fichier}

\nkfile{travail/boards/ma\_grille.json}. Le chemin exact est affiché
dans le panneau \emph{Règles}, section \textbf{Plateau}. Le fichier apparaît dans
la liste déroulante à la seconde suivante — \textbf{sans compilation}, puisqu'il
n'y a rien à compiler.

Au premier lancement, l'atelier y écrit un exemple :

\begin{quote}\small
Écrit un exemple si le dossier est vide, à partir du plateau que le moteur chargé
expose. \textbf{Le format documenté vaut mieux qu'un format décrit} : celui-ci est
forcément valide, puisqu'il vient du module.
\hfill\textsc{NkcBoardLibrary.h}
\end{quote}

\section{Le format, sur un exemple minuscule}

Un plateau carré $3\times3$, deux joueurs en diagonale, une case bloquée au
centre :

\begin{nkcode}{Ecrit pour ce cours — format de LoadBoardJson}
{
  "topology": "SQUARE_4",
  "cells": [[0,0],[1,0],[2,0],
            [0,1],[1,1],[2,1],
            [0,2],[1,2],[2,2]],
  "blocked": [[1,1]],
  "starts": [
    {"player": 0, "q": 0, "r": 0, "level": 0},
    {"player": 1, "q": 2, "r": 2, "level": 0}
  ],
  "min_players": 2,
  "max_players": 2
}
\end{nkcode}

\begin{center}
\begin{tabular}{@{}p{2.6cm}p{4.4cm}p{6.4cm}@{}}
\toprule
\textbf{Champ} & \textbf{Rôle} & \textbf{Piège} \\
\midrule
\texttt{topology} & hex ou carré & fixe le voisinage \emph{et} l'interprétation de $(q,r)$ \\
\texttt{cells} & \textbf{la forme réelle} & une coordonnée absente d'ici n'existe pas : c'est le hors-plateau \\
\texttt{blocked} & cases inoccupables & doivent aussi figurer dans \texttt{cells} \\
\texttt{starts} & totems de départ & \texttt{player} au-delà de \texttt{max\_players} est ignoré \\
\nkfile{min_players} & bornes du plateau & informatif, avec \nkfile{max_players} :
  c'est le moteur qui décide quoi en faire \\
\bottomrule
\end{tabular}
\end{center}

\begin{nkpiege}[\texttt{cells} définit la forme, pas un rectangle englobant]
C'est ce qui permet des plateaux en croix, en diamant, ou troués. Ne listez
\textbf{que} les cases qui existent. Le moteur de référence teste l'appartenance
par recherche dans cette liste : une coordonnée absente est du hors-plateau,
exactement comme si elle était au-delà du bord.
\end{nkpiege}

\section{Hexagone : la conversion qui piège}

En hexagone, les coordonnées sont \textbf{axiales} $(q, r)$, pas des lignes et des
colonnes. Le moteur de référence construit son $6\times7$ ainsi :

\begin{nkcode}{modules/rules/ConquerorRulesV2.cpp — BuildDefaultBoard}
for (int32 row = 0; row < 7; ++row) {
    for (int32 col = 0; col < 6; ++col) {
        NkcCoord c;
        c.q = static_cast<int16>(col - (row >> 1));  // odd-r -> axial
        c.r = static_cast<int16>(row);
        B.cells[B.cellCount++] = c;
    }
}
\end{nkcode}

La ligne qui compte est \texttt{col - (row >> 1)} : c'est la conversion
\textbf{odd-r vers axial}. Si vous écrivez un générateur de plateau hexagonal,
c'est là que vous vous tromperez. Le symptôme est net et reconnaissable : le
plateau s'affiche en \textbf{losange penché} au lieu d'un rectangle d'hexagones.

\begin{nkretenir}[Le raccourci qui évite l'erreur]
Ne calculez pas vos coordonnées à la main. Chargez le plateau par défaut, cliquez
\textbf{« Exporter le plateau courant »} dans le panneau \emph{Règles}, et
modifiez le fichier obtenu. Vous partez d'une géométrie juste.
\end{nkretenir}

\section{La symétrie n'est pas une coquetterie}

\begin{quote}\small
Le placement initial doit être symétrique sous une isométrie du plateau. Le moteur
de test doit refuser de lancer un batch d'équilibrage sur un plateau asymétrique,
ou signaler l'asymétrie dans le rapport. Sans cela, \textbf{tout écart de winrate
mesuré est ininterprétable} : on ne saura pas s'il vient des règles ou du plateau.
\hfill\textsc{regles\_completes\_v2.md:133}
\end{quote}

\begin{nkpiege}[Cette vérification n'existe pas encore dans l'atelier]
Le document l'exige ; l'atelier ne la fait pas. \textbf{C'est à vous de vérifier
la symétrie de vos plateaux à l'œil} avant de lancer une campagne dessus.

Il y a toutefois un garde-fou partiel : la campagne \textbf{inverse les côtés une
partie sur deux} (chapitre 6), ce qui sépare l'avantage de position de l'avantage
de stratégie. Cela compense un plateau asymétrique \emph{entre les deux camps},
mais pas un plateau qui avantagerait structurellement le premier joueur.
\end{nkpiege}

\section{Quand le moteur refuse votre fichier}

L'atelier ne parse pas le JSON : il passe la chaîne telle quelle à
\texttt{LoadBoardJson}. C'est donc \textbf{le module} qui accepte ou refuse.

\begin{nkcode}{src/ConquerorLab/NkcBoardLibrary.h — LoadInto}
if (!vt.LoadBoardJson(inst, text.CStr())) {
    mMsg  = "Le moteur a REFUSE ce plateau : ";
    mMsg += mFiles[idx].name;
    mMsg += "  (cellules manquantes ou JSON invalide)";
    return false;
}
\end{nkcode}

Ce message s'affiche en rouge dans le panneau \emph{Règles}. Trois causes, par
ordre de fréquence :

\begin{enumerate}
  \item \textbf{\texttt{cells} vide ou absent} — le moteur de référence refuse
        explicitement un plateau sans cases ;
  \item \textbf{JSON malformé} — le scanner des modules est minimal : ni
        commentaires, ni virgules traînantes ;
  \item \textbf{le module ne sait pas charger de plateau} —
        \nkfile{RegleContratNu.cpp} renvoie toujours 0 : son plateau est figé dans
        le code. C'est un renoncement assumé pour un exemple.
\end{enumerate}

\begin{nkretenir}[Conséquence heureuse de « l'atelier ne parse rien »]
Si votre moteur accepte un champ supplémentaire — disons
\texttt{"artefacts":[\{"q":1,"r":2,"id":7\}]} — il le verra \textbf{sans qu'on
touche à l'atelier}. Le format du plateau vous appartient ; le contrat ne fixe
qu'un socle commun.
\end{nkretenir}

\section{Ce que l'atelier ne fait pas}

Le document de règles parle d'un plateau « éditable à la souris ». \textbf{Cet
éditeur n'existe pas.} Aujourd'hui, on écrit le JSON à la main, ou on exporte puis
on modifie.

C'est une limite réelle et assumée : l'éditeur graphique viendra quand quelqu'un
aura assez souffert du fichier texte pour savoir de quoi il a besoin.

\begin{nkexo}[1 — Le plateau troué]
Exportez le plateau par défaut, puis retirez huit cases au centre pour en faire un
anneau. Chargez-le. Que devient la durée moyenne d'une partie sur 200 parties IA
contre IA ? Expliquez.
\end{nkexo}

\begin{nkexo}[2 — Carré contre hexagone]
Fabriquez un plateau \texttt{SQUARE\_8} de 42 cases ($6\times7$) et comparez-le à
l'hexagone $6\times7$ par défaut, \textbf{à même nombre de cases}. Lancez 500
parties sur chacun. Le taux de nuls change-t-il ? Le nombre de coups ? Que dit ce
résultat sur l'effet du voisinage ?
\end{nkexo}

\begin{nkexo}[3 — L'asymétrie visible]
Fabriquez volontairement un plateau asymétrique : trois totems de départ au joueur
0, deux au joueur 1. Lancez 400 parties \emph{avec} l'inversion des côtés, puis
400 \emph{sans}. Comparez et expliquez précisément ce que l'inversion a corrigé —
et ce qu'elle n'a pas corrigé.
\end{nkexo}


\section{La forme des cellules}

Une grille carrée doit-elle se dessiner en carrés ? Non — et les séparer est
instructif.

\begin{nkcode}{boards/\_generer.py — plateau()}
def plateau(topology, cells, starts, blocked=None, cell_shape=None):
    """`cell_shape` est de la PRESENTATION : "SQUARE", "HEX" ou "CIRCLE".

    Elle ne touche ni au voisinage, ni aux coups legaux, ni au resultat —
    c'est justement pourquoi elle n'est PAS dans le contrat : le moteur ne
    la lit jamais, seul l'atelier la regarde pour dessiner.
    """
\end{nkcode}

Trois formes livrées : \textbf{carrée}, \textbf{hexagonale}, \textbf{ronde}. Le
champ est facultatif ; absent, l'atelier déduit la forme de la topologie, comme
avant. Le panneau \emph{Règles} offre en plus un réglage \textbf{Forme des
cellules} qui prime sur le fichier — pour comparer sans rien réécrire.

\begin{nkretenir}[Trois sources, un ordre de priorité]
\begin{enumerate}
  \item le module, s'il fournit \texttt{GetCellShape} (chapitre 5) — il a le
        dernier mot sur \emph{sa} géométrie ;
  \item votre choix dans le panneau \emph{Règles} ;
  \item ce que déclare le fichier \texttt{.json}.
\end{enumerate}
\end{nkretenir}

\section{La bibliothèque livrée}

Quatorze plateaux, présents dès le premier lancement.

\begin{center}
\begin{tabular}{@{}p{3.6cm}p{10.2cm}@{}}
\toprule
\textbf{Forme} & \textbf{Fichiers} \\
\midrule
hexagonal & \texttt{hexagone\_6x7} (référence), \texttt{4x5}, \texttt{9x8}, \texttt{coeur\_bloque}, \texttt{plat} \\
rectangulaire & \texttt{rectangle\_8x6}, \texttt{carre\_8x8\_diagonales} \\
en croix & \texttt{plus} \\
parallélogramme & \texttt{parallelogramme\_6x7}, \texttt{parallelogramme\_8x5} \\
libre & \texttt{diamant} \\
formes de cellule & \texttt{rectangle\_8x6\_rond}, \texttt{hexagone\_6x7\_rond}, \texttt{carre\_8x8\_hexagones} \\
\bottomrule
\end{tabular}
\end{center}

Le parallélogramme mérite un mot : c'est \texttt{hex\_rect} \textbf{sans} la
correction \texttt{- (row >> 1)}, autrement dit la forme qu'on obtient quand on
oublie la conversion odd-r → axial. La livrer comme un plateau à part entière
vaut mieux que de la laisser apparaître comme un bug — la reconnaître à l'écran
fait comprendre du même coup ce que corrige \texttt{hex\_rect}.

Tous sont produits par \nkfile{boards/_generer.py}, jamais écrits à la main : la
symétrie des départs exigée par \textsc{regles §4.2} y est \textbf{calculée par
réflexion}. Un plateau asymétrique rend tout écart de winrate ininterprétable.
